
\documentclass[12pt,a4paper]{article}
\usepackage[utf8]{inputenc}
\usepackage[T1]{fontenc}
\usepackage[brazil]{babel}
\usepackage{amsmath,amssymb,amsthm}
\usepackage{geometry}
\geometry{margin=2.5cm}

\newtheorem{theorem}{Teorema}[section]

\title{Teorema de Euclides: existem infinitos números primos}
\author{}
\date{}

\begin{document}
\maketitle

\begin{theorem}[Euclides]
Existem infinitos números primos.
\end{theorem}

\begin{proof}
Suponhamos, por absurdo, que exista apenas uma quantidade finita de
números primos. Listemo-los como
\[
  p_1,\; p_2,\; \dots,\; p_n,
\]
todos os primos existentes segundo essa hipótese.

Considere o número
\[
  N = p_1 p_2 \cdots p_n + 1.
\]
Como $N \geq 2$, pelo Teorema Fundamental da Aritmética $N$ possui ao
menos um divisor primo $q$. Pela hipótese, $q = p_i$ para algum
$i \in \{1,\dots,n\}$.

Mas $q \mid N$ e $q \mid p_1 p_2 \cdots p_n$ acarretam
\[
  q \mid \bigl(N - p_1 p_2 \cdots p_n\bigr) = 1,
\]
o que é impossível, pois $q \geq 2$.

Obtivemos uma contradição. Logo a suposição é falsa e existem
infinitos números primos.
\end{proof}

\end{document}
